\documentclass[letterpaper]{article}
\usepackage[submission]{aaai2027}
\usepackage[hyphens]{url}
\usepackage{graphicx}
\urlstyle{rm}
\def\UrlFont{\rm}
\usepackage{natbib}
\usepackage{bibentry}
\usepackage{caption}
\frenchspacing
\usepackage{booktabs}
\usepackage{amsmath}
\usepackage{amssymb}
\usepackage{listings}
\newcommand{\op}{\diamond}
\setcounter{secnumdepth}{2}

% Verbatim prompt blocks: map the non-ASCII characters that appear in the
% prompt text so pdfLaTeX does not choke inside lstlisting.
\lstset{%
  basicstyle={\footnotesize\ttfamily},
  breaklines=true, showstringspaces=false, frame=single,
  xleftmargin=0pt, columns=fullflexible,
  literate={◇}{{$\diamond$}}1 {—}{{---}}1 {–}{{--}}1
           {•}{{$\bullet$}}1 {→}{{$\rightarrow$}}1 {✓}{{$\checkmark$}}1
}

\pdfinfo{/TemplateVersion (2027.1)}

\begin{document}

\appendix

\renewcommand{\thetable}{A\arabic{table}}
\renewcommand{\thefigure}{A\arabic{figure}}
\setcounter{table}{0}
\setcounter{figure}{0}

\section*{Appendix}


\section{Worked Example: Laws 12857 and 33436}
\label{app:worked}

References to sections, tables and figures without an A-prefix refer to the
main paper, and citations resolve in the main paper's reference list.

Among the laws the portfolio resolves is a pair the Equational Theories
Project \citep{bolan2025equational} recorded as having only trivial finite
models, with the existence of an infinite model left unresolved. In the ETP numbering they are laws 12857 and
33436:
\begin{align*}
L_{12857}\colon\quad & x = y \op ((x \op (y \op (z \op z))) \op y), \\
L_{33436}\colon\quad & x = ((y \op (((z \op z) \op y) \op x)) \op y).
\end{align*}
The two are duals: reversing the arguments of every occurrence of $\op$ in one
yields the other. A magma satisfies one exactly when its opposite magma
satisfies the other, so the pair is a single open case up to duality, and
resolving either resolves both.

\subsection{Two independent certificates}

Both laws were resolved twice, by methods that share no search strategy.

Vampire 5.0.1 \citep{kovacs2013first} was run in ordered saturation mode (\texttt{-sa otter}, default
Knuth--Bendix ordering) on the law together with the assertion that two
elements are distinct. It terminated with \texttt{SZS status Satisfiable} in
7.95 seconds and 126\,MB of memory, printing an active set of 357 unit
equations. A terminating, complete saturation containing no contradiction
describes a model, as set out in Section~3.3; here the model is one in which
$x = y$ fails, and the admissibility certificate the instance already carries
forces that model to be infinite.

Twee 2.6.1 \citep{smallbone2021twee} was run independently on both laws and returned \texttt{SZS status
CounterSatisfiable} for each, deriving 55 rules for 12857 and 54 for 33436.
Unfailing completion is complete by construction, so this is a second
certificate for the same conclusion by a different route.

\subsection{Verifying the presentation}

The saturated presentation is checked directly rather than taken on the
prover's word. Three properties are confirmed. The ordering check verifies that
the ordering used to evaluate the presentation matches the one recorded in the
certificate header, which is necessary because ground confluence is inherited
from the saturation only with respect to the ordering the prover saturated
under; every pre-ordered equation is comparable under that ordering. The law
then holds on every ground instance tested, 27 of 27, and rewrites fire in all
of them, so the model is not satisfying the law vacuously. Finally the two
Skolem constants of the distinctness axiom normalise to different terms, which
establishes that the carrier has at least two elements.

The presentation is too large to reproduce here; the certificate ships in the
supplementary material, and \texttt{ordered\_model.py} reproduces all three
checks from it.

\subsection{Why this model is not Lean-checkable today}

The presentation is not an off-the-shelf term rewriting system, which is why
this model is graded by the prover certificate rather than by the Lean kernel
\citep{moura2021lean}.
Of the 357 equations, a substantial minority carry a free variable on each
side, so they cannot be oriented into rewrite rules in the usual sense.

That is a constraint on printing the system, not on the construction. The model
rewrites a ground term whenever the ground instance strictly decreases, and a
variable unbound on the other side of an equation may be mapped to any ground
term that makes the instance decrease --- in practice the smallest constant.
Every equation is therefore usable, in whichever direction decreases the
instance. Termination is immediate because the ordering is well founded on
ground terms, and ground confluence is inherited from saturation under
unfailing completion.

What is missing is an off-the-shelf certifier. Tools such as CSI and TTT2 check
plain rewriting systems, and the object here is an ordered one. A Lean
formalization of this model therefore awaits a machine-checked ground
confluence proof. Models of the algebraic kind, by contrast, are Lean-checkable
today; no admissible law admits one, for the reason given in Section~3.3.


\section{LLM Evaluation Protocol and Prompts}
\label{app:protocol}

\subsection{The fixed protocol}

Every solver receives the same three inputs: the law, the task statement, and a
fixed answer format. A submission may be sent to the judge up to three times,
and each verdict is returned verbatim as feedback before the next attempt. No
solver receives a hint that is not derived mechanically from the law itself,
and the harness is frozen across all language models, so differences in outcome
are attributable to the model rather than to the scaffolding.

Saturation provers are withheld from the language-model setting. With access to
them a solver could re-run the automated portfolio of Section~4.1 rather than
construct anything; the judge and the autoformalization harness are the tools
available to any solver built on ALPS.

All LLM calls use provider-default sampling through a single API gateway. The
APIs expose no seed parameter, and the reasoning LLMs are nondeterministic at
any setting, so bitwise reproduction of a run is not available. Run-to-run
variation is therefore measured rather than controlled: the reproduction run
of Section~4.2 repeats the nine solved and five matched unsolved laws under
the identical frozen harness. The result files in the supplementary material
record every attempt; the submitted chain is withheld from solved rows, since
a verified chain is a solution to a certified-easy law, and every reported
count is reproducible from the retained fields.

\subsection{The trivial channel}

On the trivial side the formalization burden is removed. The model emits a
chain of equalities between two arbitrary carrier elements $a$ and $b$, each
step being one application of the law, and returns it as JSON. The harness
matches each step against the law, bridges gaps of at most three single
applications automatically, and assembles the Lean proof itself. The model
therefore searches for a derivation rather than writing Lean.

A step is accepted only when some substitution instance of one side of the law
rewrites one term into the next at exactly one position, with everything
outside that position unchanged. Two rejections are absolute: relating $a$ and
$b$ directly, and any jump the bridging search cannot close.

A \emph{waypoint lemma} is an equality a prover has already derived from the
law, supplied in the prompt. Waypoints shorten planning --- the model can build
a chain to a waypoint and continue from there --- but they cannot repair an
individual step, since every step the model writes must still be a correct
application of the law. Comparing runs with and without waypoints therefore
separates failures of chain-finding from failures of step validity.

\subsection{The construction channel}

On the construction side the model proposes a finite presentation $E$: a set of
unconditional equations over $\op$, returned as JSON. The submission format
admits only unconditional equations, so implications, disequations and literal
restatements of $x = y$ never reach the prover. The judge then runs the two
checks of Equation~(6).

\subsection{Prompts}

Figure~\ref{fig:prompt-trivial} gives the trivial-channel prompt and
Figure~\ref{fig:prompt-construct} the construction-channel prompt, both
verbatim. Braces marked \texttt{\{law\}} and \texttt{\{hint\_block\}} are
substituted at call time; the hint block is empty in the no-waypoint condition
and lists the prover-derived waypoint lemmas otherwise. On a rejected
submission the harness appends the judge's reason to the next round's prompt:
\texttt{Your previous chain was REJECTED: \{feedback\}} on the trivial side,
and \texttt{Your previous proposal failed to certify: \{feedback\}} followed by
\texttt{Try a different presentation.} on the construction side.

\begin{figure*}[t]
\begin{lstlisting}
You are proving a magma law forces triviality. You will NOT write Lean — you
give only a sequence of terms, and a theorem prover checks each step.

Law:  {law}
Read it as a rewrite rule. Write the law as  L = R, where L is the single
variable on the left and R is the big right-hand term. ONE step does exactly
one of:
  • pick any subterm `e` of the current term and replace that ONE occurrence
    with R, taking L := e (so `e` becomes R with the law's left variable set
    to `e`); OR
  • the reverse — replace a subterm that exactly matches an instance of R by
    the corresponding `e`.
Everything OUTSIDE that one chosen subterm stays byte-for-byte identical.

Goal: build  a = t1 = t2 = ... = b. Adjacent terms should be only a FEW (1–3)
such steps apart — you may skip intermediate terms and the checker will fill
short gaps automatically, so give WAYPOINTS, not necessarily every atomic
step. Prefer smaller jumps.

HARD RULES — a jump breaking any of these is rejected, so self-check each one
before writing it:
  • You may NOT write "a = b" or otherwise relate a and b directly. They are
    opaque atoms; the only way to connect them is through the law.
  • Each jump must be a GENUINE short derivation — at most ~3 single law
    applications apart. A leap between unrelated terms cannot be bridged and
    is rejected.
  • Every application is a REAL instance: a subterm `e → R[L:=e]` (or its
    reverse) for some choice of the law's other variables. If you can't name
    that instance, the step is invalid.

STRATEGY: rewrite FORWARD from `a` (each rewrite grows the term via the law)
toward a term the law forces to collapse — these laws let you derive
`op(_, _) = (any element)`; once both `a` and `b` reduce to a common term, the
chain closes.{hint_block}

Return ONLY JSON:  {"chain": ["a", "<term>", ..., "b"]}   first "a", last "b";
use only a, b, ◇, ().
Example (toy law  x = y ◇ y):  {"chain": ["a", "(a ◇ a)", "b"]}
  step 1: `a` → `(a ◇ a)` — the law with x:=a, y:=a gives  a = a◇a. ✓
  step 2: `(a ◇ a)` → `b` — reverse: the law with x:=b, y:=a gives  b = a◇a,
    so  a◇a = b. ✓
Only write steps you can justify this way.
\end{lstlisting}
\caption{The trivial-channel prompt, verbatim. The model returns a chain of
terms as JSON; the harness matches each step against the law, bridges gaps of
at most three single applications, and assembles the Lean proof.}
\label{fig:prompt-trivial}
\end{figure*}

\begin{figure*}[t]
\begin{lstlisting}
You are constructing an infinite model for a magma law, to be checked by an
automated theorem prover.

The law is:  {law}
(the binary magma operation is written with the diamond; the law holds for all
inputs.)

This law has NO nontrivial FINITE model, so its model is necessarily infinite,
and it CANNOT be an arithmetic formula (any polynomial op over the integers
would descend to a finite model). Instead, propose a finite set of DEFINING
EQUATIONS for the operation — a presentation of the model — from which the law
follows and which is consistent with the carrier having two distinct elements.

Think of it as the completed rewrite system a Knuth-Bendix procedure would
converge to: a handful of equations that (a) entail the law, and (b) do not
force all elements equal.

Return ONLY a JSON object, nothing else:
  {"E": ["<equation>", "<equation>", ...]}
Each equation uses variables x,y,z,w and the diamond operator, in the form
LHS = RHS .
Example shape (NOT a solution):
  {"E": ["x ◇ (y ◇ x) = y", "(x ◇ x) ◇ x = x ◇ x"]}
Do NOT just restate the law, and do NOT write `x = y`.
\end{lstlisting}
\caption{The construction-channel prompt, verbatim. The model returns a finite
presentation $E$ as JSON, which the judge submits to the two checks of
Equation~(6).}
\label{fig:prompt-construct}
\end{figure*}

% ============================================================================
\section{Automated Baseline: Per-Configuration Results}
\label{app:perconfig}

The portfolio spans eight configurations: Vampire 5.0.1 in five modes, E 3.3.5
\citep{schulz2019faster} in two, and Twee 2.6.1. Table~\ref{tab:perconfig} reports how many distinct
laws each configuration resolves at each per-configuration budget, over the
4,141-law residual. The counts overlap, since several configurations resolve
the same law, so the column totals exceed the 114 laws resolved overall.

\begin{table*}[t]
\centering
\begin{tabular}{@{}lrrrrrrr@{}}
\toprule
Configuration & 30 & 60 & 120 & 300 & 600 & Total & Models \\
\midrule
twee/complete         & 69 & 5 & 4 & 3 & 5 & 86 & 4 \\
eprover/triv/auto     & 64 & 2 & 2 & 1 & 1 & 70 & 0 \\
vampire/triv/casc+lpo & 30 & 1 & 1 & 1 & 0 & 33 & 0 \\
vampire/triv/casc     & 24 & 1 & 0 & 1 & 0 & 26 & 0 \\
vampire/triv/discount & 22 & 0 & 0 & 0 & 0 & 22 & 0 \\
vampire/sat/otter+lpo & 13 & 1 & 0 & 0 & 0 & 14 & 0 \\
vampire/sat/otter+kbo & 12 & 1 & 1 & 0 & 0 & 14 & 0 \\
eprover/sat/satauto   &  0 & 0 & 0 & 0 & 0 &  0 & 0 \\
\bottomrule
\end{tabular}
\caption{Distinct laws resolved by each configuration at each per-configuration
budget in seconds, over the 4,141-law residual. Counts overlap across
configurations. ``Models'' counts laws resolved by exhibiting a nontrivial
model rather than by proving triviality.}
\label{tab:perconfig}
\end{table*}

Two features of the table bear on the main paper's claims. First, every
configuration concentrates its resolutions at the 30-second rung; no
configuration shows a profile in which added budget yields resolutions at a
steady rate. Second, all four models the submitted sweep recovers come from a
single configuration, \texttt{twee/complete}. The Vampire saturation modes
resolve 14 laws each, all by proving triviality. The construction side of the
portfolio therefore rested on one configuration.

The sweep exits a law's ladder as soon as any configuration resolves it, so the
number of laws reaching each rung declines: 4,141 at 30 seconds, then 4,050,
4,042, 4,037 and 4,033.

\section{LLM Failure Analysis}
\label{app:failures}

\begin{table}[t]
\centering
\begin{tabular}{@{}lrrrr@{}}
\toprule
LLM & Too strong & Too weak & Both fail & One check \\
\midrule
o3      & 23 & 2  & 0  & 25 \\
o4-mini & 11 & 4  & 10 & 15 \\
GPT-4.1 & 10 & 11 & 4  & 21 \\
\bottomrule
\end{tabular}
\caption{Final-round classification of all 25 hard-tier construction
submissions per LLM. ``Too strong'' entails the law but admits only
one-element models (passes check (a), fails (b)); ``too weak'' maintains two
distinct elements but fails to entail the law (passes (b), fails (a)). The last
column counts laws on which the LLM clears exactly one of the two checks.}
\label{tab:failures}
\end{table}

Each language model attempts all 25 hard-tier laws over three feedback rounds,
and no submission from any model passes both checks. Because a run ends either
in acceptance or after three rejections, every law leaves a final-round
submission, and Table~\ref{tab:failures} classifies all 75 of them by which of
the two checks that final submission satisfies.

The three classes partition the outcomes. A presentation is \emph{too strong}
when it entails the law but admits only one-element models: check~(a) succeeds
and check~(b) fails. It is \emph{too weak} when it keeps two elements distinct
but does not entail the law: check~(b) succeeds and check~(a) fails. A
submission failing both satisfies neither condition. The final column, counting
laws on which the model clears exactly one check, is therefore 25 minus the
both-fail count.

The trivial-side failures follow one dominant pattern at every setting, and it
is a different kind of failure. The submitted chains contain steps that are not
applications of the law: no substitution instance of either side rewrites one
term into the next, so the harness rejects the chain at the matching stage
rather than at the prover. The failure is in executing individual steps
soundly, not in planning a route between $a$ and $b$.

\section{Corpus Yield and Renewability}
\label{app:renew}

\begin{table}[t]
\centering
\begin{tabular}{@{}rrrr@{}}
\toprule
Order & Screened & Austin & Austin per 1{,}000 \\
\midrule
5 & 115   & 11  & 95.7$^{\ddagger}$ \\
6 & 1,418 & 29  & 20.5 \\
7 & 4,803 & 117 & 24.4 \\
8 & 4,138 & 105 & 25.4 \\
\bottomrule
\end{tabular}
\caption{Admissible-law yield by order, over all 10,474 screened candidates.
Density rises monotonically across the orders the extension engine generates,
so the supply does not thin as the corpus grows. $^{\ddagger}$Order 5 is
excluded from the trend: those laws come from ETP's curated enumeration rather
than a uniform draw from the generator.}
\label{tab:renewability}
\end{table}

Section~3.2 claims that the corpus has no terminal order and that fresh
instances can therefore be minted after any training cutoff. That claim is
worth checking against yield: a generator whose output thins as the order
grows would eventually stop supplying instances, whatever the recursion
permits in principle.

Table~\ref{tab:renewability} gives the yield of the screening pipeline broken
down by the order of the candidate law, over all 10,474 candidates screened.

Yield must be read as a density rather than a count. The raw number of Austin
laws dips from 117 at order 7 to 105 at order 8, but fewer order-8 candidates
were screened, so the dip measures how much was generated rather than how much
was there. Normalised to admissible laws per thousand screened candidates, the
yield rises across every order the extension engine produces: 20.5 at order 6,
24.4 at order 7, and 25.4 at order 8. Order 5 is excluded from the comparison
because those laws are ETP's curated enumeration rather than a uniform draw
from the generator, which is why its density is several times higher.

Counting only laws that open a new equivalence class gives 16.9, 19.4 and 16.2
per thousand across the same three orders. That density does not fall, so the
corpus continues to produce genuinely distinct problems rather than variations
on classes it has already covered; unlike the overall yield, it is not
monotonic, and we do not claim it is rising.


\section{Validation of the Automated Judge}
\label{app:judge}

Both channels of the judge run a control suite before any evaluation. The
suites are part of the released code and can be re-run with
\texttt{-{}-selftest}.

\paragraph{The trivial channel.} A submission is a Lean proof of a proposition
the judge generates, and two things must hold: the proof must elaborate at
exactly the generated type, and its axiom footprint must lie within the
allowlist $\{$\texttt{propext}, \texttt{Quot.sound},
\texttt{Classical.choice}$\}$. Anything outside that set is a failure, and
\texttt{sorryAx}, \texttt{Lean.ofReduceBool} and \texttt{Lean.trustCompiler}
are named explicitly as forbidden. The footprint is an allowlist rather than a
blocklist so that an unanticipated axiom fails closed.

Seven negative controls exercise the ways a submitter could avoid proving
anything, and each is rejected by a textual scan that runs before Lean is
invoked, since several of them would otherwise subvert Lean itself: a
\texttt{sorry}; a declared axiom asserting the goal; two forms of namespace
shadowing; redefining the law predicate as \texttt{True}; \texttt{native\_decide};
and naming the submitted theorem anything other than the identifier the judged
file elaborates. A positive control confirms the suite is not simply rejecting
everything: a reference answer passes the scan and compiles with an allowed
footprint.

\paragraph{The construction channel.} Three controls cover the three ways a
presentation can fail to encode a construction, matching the cases discussed in
Section~3.3. Submitting the law itself serves as the positive control: it
certifies exactly when bare saturation of $L \cup \{\,a \neq b\,\}$
terminates, which screening has already established for the seed law used
here. On the hard tier the same submission diverges at check~(b), so echoing
the law back does not produce a certificate there. The left projection $\{x \op y = x\}$ is too weak: it admits models
with distinct elements but does not entail the law, so check~(a) fails. The
collapse $\{x = y\}$ is too strong: it entails the law but admits only
one-element models, so check~(b) fails. The suite passes only when the first
certifies and the other two do not.


\section{Diversity of the Constructed Models}
\label{app:diversity}

Section~3.2 reports corpus size in equivalence classes, so that laws posing the
same problem are not counted twice. Logical equivalence is the natural
criterion, but it leaves a question open: two laws that are not equivalent may
still be satisfied by essentially the same construction, in which case the
class count would overstate how many distinct problems the corpus contains.

We test this directly, without a prover, by asking which laws each constructed
model satisfies. Every model built during screening is evaluated against every
one of the 262 laws proved Austin, giving 68,382 ordered pairs and no missing
models. A model transfers to a law when it satisfies that law as well as its
own.

The result is that construction classes and logical classes coincide exactly:
both number 195. The pairs on which transfer runs in both directions number
255, precisely the 255 pairs the prover proved equivalent, so mutual
model-satisfaction and logical equivalence pick out the same relation --- a
cross-check on the deduplication of Section~3.2 by a method that shares none of
its machinery. Transfer across class boundaries is rare and one-directional:
173 of 68,382 ordered pairs, or 0.25\%. A typical model satisfies only the law
it was built for, with a median coverage of one law, a mean of 3.67 and a
maximum of 19. No cell is undecided, and every model satisfies its own law.

The 195 classes are therefore construction-distinct and not merely logically
distinct. The models do not collapse onto a shared construction that a solver
could learn once and reapply.


\section{Computing Infrastructure and Cost}
\label{app:infra}

All prover experiments ran on a dual-socket AMD EPYC 7313 machine --- 2 sockets
$\times$ 16 cores with SMT disabled, giving 32 physical cores --- with 503\,GB
of RAM under Ubuntu 22.04.5 LTS. The portfolio sweep was sharded 32 ways, one
shard per physical core, so no configuration competed with a sibling
hyperthread for its time budget. Language-model experiments were run through
the OpenRouter API.

Table~\ref{tab:cost} reports token usage for every language-model run described
in the paper. The two certified-easy runs account for more than half of the
total, because each of the 63 laws may consume up to three feedback rounds and
because reasoning traces on that set are long: a solved law costs a median of
49,221 completion tokens. The construction runs are cheaper per law despite
resolving nothing, since a presentation is a short object to emit. GPT-4.1's
totals are an order of magnitude below the reasoning models' at comparable call
counts, which reflects the absence of a reasoning trace rather than a shorter
answer.

Files with more than 63 records contain repeated passes under identical
settings, appended in run order. The no-waypoint o3 file holds a first pass
over all 63 laws, on which four chains certify, followed by a second pass over
the 36 laws then unsolved, which adds two more; the six no-waypoint solutions
reported in the main text are the union of the two passes. The waypoint o3
file holds a first pass over all 63 laws, on which nine chains certify,
followed by two further passes over 32 of the unsolved laws, which add none.
The reproduction run is a separate file.

\begin{table}[t]
\centering
\begin{tabular}{@{}lrrr@{}}
\toprule
Run & Records & Completion & Prompt \\
\midrule
o3 cert-63, waypoints    & 127 & 3,094,531 & 164,044 \\
o3 cert-63, no waypoints & 99  & 3,289,573 & 130,963 \\
o3 low effort            & 9   & 152,309   & 20,724 \\
o3 reproduction          & 14  & 575,997   & 59,653 \\
o4-mini cert-63          & 63  & 974,972   & 147,300 \\
GPT-4.1 cert-63, waypoints    & 63 & 61,670 & 149,538 \\
GPT-4.1 cert-63, no waypoints & 63 & 71,736 & 139,620 \\
GPT-4.1 early run (95 laws)   & 100 & 192,466 & 269,443 \\
GPT-4.1 pilot (10 laws)  & 10  & 9,234     & 31,870 \\
o3 trivial, hard-25      & 25  & 1,226,978 & 53,475 \\
o3 construct, hard-25    & 25  & 1,120,628 & 144,122 \\
o4-mini construct        & 25  & 554,835   & 24,818 \\
GPT-4.1 construct        & 46  & 6,605     & 45,755 \\
\midrule
Total                    & 669 & 11,331,534 & 1,381,325 \\
\bottomrule
\end{tabular}
\caption{Token usage across all reported language-model runs. A solved
certified-easy law costs a median of 49,221 completion tokens.}
\label{tab:cost}
\end{table}

% Resolve in-text citations against the main paper's bibliography without
% printing a reference list (all cited keys appear in the main paper).
\nobibliography{aaai2027}

\end{document}